\begin{textsample}{NEG Dim 4 – global\_south\_2024\_04 – Score 1.00 – global\_south\_2024\_04/107836564.txt}  \label{ex:f4_neg_196}
The sound of solidarity One of the numerous war crimes that South Africa accused Israel of in its case of genocide before the International Court of Justice was its obliteration of Gaza ' s cultural heritage."Israel has damaged and destroyed numerous centres of Palestinian learning and culture,"South Africa stated at the Peace Palace in The Hague in January.
By that time, nearly 200 sites of historical importance had been destroyed or damaged in Israeli air raids on Gaza, according to an Al Jazeera article on 14 January.
By now there must be several more such sites of destruction.
As scholar Caitlin Procter wrote in Jacobin,"The destruction of ideas, creativity, and heritage in Gaza, in this landscape already harrowed by siege and occupation, can not be depicted as collateral damage in a supposed war against Hamas."Alongside the horror of mass killing, it is vital not to ignore this deliberate destruction of the cultural life that makes Gaza such a rich tatreez a form of traditional Palestinian embroidery of intellect, beauty, and @ @ @ @ @ @ @ @ @ @ been expressing solidarity with war-torn Gaza and its beleaguered people.
Many, including here in South Africa, have created art works, music and writing to show their support for Gazans, especially its artists and cultural workers.
Last month, the South African art and music collective, Gusher Sound, in conjunction with African Artists Against Apartheid, released a 10-track album titled Piercing the Silence on the online music distribution platform, Bandcamp.
The artists are based in either Johannesburg or Cape Town, apart from UK musician, Glor1A, and Lover Buffet, who is from Berlin.
Gusher ' s two young leading lights, managing director Lunathi Kondlo and creative director Mike Ngudle, are both politically aware.
In December they were watching the horrors unfold in Gaza and decided they should use their skills to respond, they told me in a Zoom interview this week."We don't operate in a void,"Ngudle said."We didn't feel right to just go about our regular schedule without acknowledging @ @ @ @ @ @ @ @ @ @ sense as a project that we could put together."They know several artists who are politically inclined and Piercing the Silence was the result when they called on them to contribute."The project is basically our contribution to the causes that we believe in,"Kondlo said.
The fundraising through the album goes beyond just Palestine, with proceeds also going to organisations providing humanitarian aid in Congo and Sudan.
Piercing the Silence is more than just a worthy project -- it is a wonderfully eclectic, gentle, laid-back and top-quality, genre-fluid compilation with folk, electronica, alt-RnB, jazz, roots, Afrobeat and beyond, from a range of seriously good artists including Kujenga, Sonder the Africanime, Thammy Ndlovu, The Mavrix, Desire Marea and Farrah Khataza.
Interestingly enough, it is not an overtly political album, with only about a third of the 10 songs focusing specifically on Palestine."Our thing was mobilisation,"said Kondlo."It was kind of music @ @ @ @ @ @ @ @ @ @ had that sense of urgency, right the idea of piercing the silence...
It might not be overtly political lyrically but there ' s definitely a sense of urgency in each song."Founded in 2020, Gusher is a multi-platform outlet for young musicians, filmmakers, designers and artists.
They will be using these skills and contacts for the next step of piercing the silence around Palestine, Congo and Sudan.
They are going to collaborate with a Johannesburg outfit called Friendly SA to make merchandise such as T-shirts and caps."We realised the relationship that people have with music and it kind of hit us that, like, people don't buy music,"explained Ngudle."So, the merchandise idea was one of the things which is how we will operate."In the meantime, Piercing the Silence is still available on Bandcamp -- proof that solidarity comes in different shapes, colours and sounds.

% matched lemmas: 
\end{textsample}
